\chapter{Аналитическая часть}

В данном разделе рассматриваются алгоритмы умножения матриц: классический алгоритм и алгоритм Винограда. Основное внимание уделяется их вычислительной сложности и особенностям реализации.

\section{Классический алгоритм умножения матриц}

Пусть имеются две матрицы: матрица $A$ размером $n \times m$ и матрица $B$ размером $m \times q$. Их произведение $C = AB$ представляет собой матрицу размером $n \times q$, элементы которой определяются следующим выражением~(\ref{eq:classic_formula}):

\begin{equation}
	c_{ij} = \sum_{k=1}^{m} a_{ik} \cdot b_{kj}, \quad i = 1,\ldots,n; \ j = 1,\ldots,q
	\label{eq:classic_formula}
\end{equation}

Сложность данного алгоритма оценивается как $O(n \cdot m \cdot q)$ операций. В случае квадратных матриц порядка $n$ сложность составляет $O(n^3)$.

\section{Алгоритм Винограда}

Алгоритм Винограда основан на альтернативном способе вычисления скалярного произведения. Для матриц $A \in \mathbb{R}^{n \times m}$ и $B \in \mathbb{R}^{m \times q}$ элемент результирующей матрицы вычисляется по формуле~(\ref{eq:winograd_main}):

\begin{equation}
	c_{ij} = \sum_{k=1}^{m/2} (a_{i,2k-1} + b_{2k,j})(a_{i,2k} + b_{2k-1,j}) - R_i - S_j
	\label{eq:winograd_main}
\end{equation}

где предварительно вычисляются вспомогательные величины~(\ref{eq:row_factor} и~\ref{eq:col_factor}):

\begin{equation}
	R_i = \sum_{k=1}^{m/2} a_{i,2k-1}a_{i,2k}, \quad i = 1,\ldots,n
	\label{eq:row_factor}
\end{equation}

\begin{equation}
	S_j = \sum_{k=1}^{m/2} b_{2k-1,j}b_{2k,j}, \quad j = 1,\ldots,q
	\label{eq:col_factor}
\end{equation}

\subsection{Анализ вычислительной сложности}

Количество операций в алгоритме Винограда:
\begin{itemize}
	\item умножения: $n \cdot \frac{m}{2} + q \cdot \frac{m}{2} + n \cdot q \cdot \frac{m}{2} = \frac{m}{2}(n + q + n \cdot q)$
	\item сложения: $n \cdot (\frac{m}{2}-1) + q \cdot (\frac{m}{2}-1) + n \cdot q \cdot (2 \cdot \frac{m}{2} + 1) = n(\frac{m}{2}-1) + q(\frac{m}{2}-1) + n \cdot q \cdot (m + 1)$
\end{itemize}

Для квадратных матриц размера $n \times n$ ($m = n$, $q = n$):
\begin{itemize}
	\item умножения: $\frac{n}{2}(n + n + n^2) = \frac{n^3}{2} + n^2$
	\item сложения: $n(\frac{n}{2}-1) + n(\frac{n}{2}-1) + n^2(n + 1) = \frac{3}{2}n^3 + \frac{3}{2}n^2 - 2n$
\end{itemize}

Таким образом, алгоритм Винограда сокращает количество умножений примерно в 2 раза по сравнению с классическим алгоритмом, но увеличивает количество сложений.

\subsection{Особенности реализации}



При нечетном значении $m$ требуется дополнительная обработка~(\ref{eq:winograd_odd}):

\begin{equation}
	c_{ij} = \left[\sum_{k=1}^{\lfloor m/2 \rfloor} (a_{i,2k-1} + b_{2k,j})(a_{i,2k} + b_{2k-1,j}) - R_i - S_j\right] + a_{i,m} \cdot b_{m,j}
	\label{eq:winograd_odd}
\end{equation}

где $R_i$ и $S_j$ вычисляются для $k = 1,\ldots,\lfloor m/2 \rfloor$.
\section*{Вывод}

Алгоритм Винограда демонстрирует выигрыш в количестве операций умножения по сравнению с классическим алгоритмом, однако требует дополнительных операций сложения и предварительных вычислений. Эффективность алгоритма Винограда проявляется на больших матрицах, где сокращение количества умножений компенсирует увеличение сложности за счет дополнительных операций.

\clearpage